\documentclass[a4paper,10pt,twoside]{article}
\usepackage{fontspec}
\setmainfont{Liberation Sans}
\setsansfont{Liberation Sans}
\renewcommand{\familydefault}{\sfdefault}
\usepackage{amsmath,amssymb}
\usepackage{geometry}
\usepackage{multicol}
\usepackage{graphicx}
\usepackage{fancyhdr}
\usepackage{lastpage}
\usepackage{booktabs}
\usepackage{listings}
\usepackage{eso-pic}
\usepackage{tikz}
\usepackage{array}
\usepackage{tabularx}
\usepackage{enumitem}
\usepackage{microtype}

\geometry{
    a4paper,
    top=1.8cm,
    bottom=1.8cm,
    left=1.2cm,
    right=1.2cm,
    headsep=14pt,
    footskip=20pt
}

\setlength{\columnsep}{16pt}
\setlength{\columnseprule}{0.35pt}

% Persistent Subtle Watermark on Every Page
\newsavebox{\watermarkbox}
\sbox{\watermarkbox}{%
    \begin{tabular}{c}
        \includegraphics[width=6.0cm]{/home/johan/GATE/GATE_CS_PYQ_DB/assets/studiva-butterfly-icon.jpg} \\[-0.2em]
        {\fontsize{14}{17}\selectfont\bfseries Studiva}
    \end{tabular}%
}
\AddToShipoutPictureBG{%
    \begin{tikzpicture}[remember picture, overlay]
        \node[opacity=0.035] at (current page.center) {\usebox{\watermarkbox}};
    \end{tikzpicture}%
}

% Running Headers & Footers
\pagestyle{fancy}
\fancyhf{}
\fancyhead[C]{\small\textbf{Studiva GATE Question Bank}}
\fancyhead[LE]{\small\textit{Computer Science \& IT (2011--2025)}}
\fancyhead[RO]{\small\textit{Studiva}}
\fancyfoot[LE,RO]{\scriptsize\textbf{Studiva} $\cdot$ \texttt{www.studiva.co.in}}
\fancyfoot[C]{\small\textbf{Page \thepage\ of \pageref{LastPage}}}
\fancyfoot[RE,LO]{\scriptsize\texttt{contact@studiva.co.in}}
\renewcommand{\headrulewidth}{0.4pt}
\renewcommand{\footrulewidth}{0.4pt}

% Listings Style (Clean, professional)
\lstset{
    basicstyle=\scriptsize\ttfamily,
    numbers=none,
    frame=single,
    rulecolor=\color{black!30},
    tabsize=2,
    breaklines=true,
    showstringspaces=false
}

\begin{document}

% ==================== PAGE 1: FORMAL TITLE & COVER PAGE ====================
\thispagestyle{empty}
\vspace*{1.5cm}
\begin{center}
    \includegraphics[width=2.5cm]{/home/johan/GATE/GATE_CS_PYQ_DB/assets/studiva-butterfly-icon.jpg}\\[1.5em]
    {\fontsize{24}{28}\selectfont\textbf{GATE COMPUTER SCIENCE \& IT}}\\[0.5em]
    {\fontsize{16}{20}\selectfont\textbf{PREVIOUS YEARS QUESTION BANK (2011 -- 2025)}}\\[0.6em]
    {\normalsize Detailed Examination Papers with Complete Official Answer Keys}\\[2.0em]
    \rule{0.65\textwidth}{0.8pt}\\[2.0em]
    {\large\textbf{Studiva GATE Examination Series}}\\[0.4em]
    {\small Comprehensive Authentic Question Bank for GATE Computer Science \& Information Technology}\\[2.8cm]

    \begin{tabular}{|ll|}
        \hline
        \textbf{Subject:} & Computer Science and Information Technology (CS/IT) \\
        \textbf{Period Covered:} & 2011 to 2025 \\
        \textbf{Examination Papers:} & 1 Complete Verified Paper Sets \\
        \textbf{Contents:} & Examination Questions \& Official Final Answer Keys \\
        \textbf{Published by:} & Studiva (\texttt{www.studiva.co.in}) \\
        \hline
    \end{tabular}

    \vfill
    {\small\textbf{Studiva $\cdot$ \texttt{www.studiva.co.in} $\cdot$ \texttt{contact@studiva.co.in}}}
\end{center}

\newpage

% ==================== PAGE 2: PUBLICATION IMPRINT & NOTES ====================
\thispagestyle{empty}
\vspace*{0.6cm}
\begin{center}
    {\large\textbf{PUBLICATION INFORMATION \& COPYRIGHT NOTICE}}\\[0.4em]
    \rule{0.6\textwidth}{0.5pt}
\end{center}

\vspace{1.2em}
\noindent\begin{tabular}{ll}
\textbf{Title:} & Studiva GATE Question Bank: Computer Science \& IT (2011--2025) \\
\textbf{Publisher:} & Studiva \\
\textbf{Website:} & \texttt{www.studiva.co.in} \\
\textbf{Official Inquiries:} & \texttt{contact@studiva.co.in} \\
\textbf{Edition:} & First Edition (2026) \\
\textbf{Syllabus Scope:} & Official GATE Syllabus (General Aptitude and CS / IT) \\
\end{tabular}

\vspace{1.8em}
\noindent\fbox{\parbox{\dimexpr\linewidth-2\fboxsep-2\fboxrule}{%
\textbf{COPYRIGHT NOTICE}\\[0.4em]
Copyright \copyright\ 2026 Studiva. All Rights Reserved.\\[0.4em]
No part of this publication may be reproduced, stored in a retrieval system, or transmitted in any form or by any means, electronic, mechanical, photocopying, recording, or otherwise, without prior written permission of Studiva, except for individual academic study.
}}

\vspace{1.5em}
\noindent\fbox{\parbox{\dimexpr\linewidth-2\fboxsep-2\fboxrule}{%
\textbf{PUBLISHER'S NOTE: ARCHIVAL QUESTION FORMAT}\\[0.4em]
This compendium has been meticulously curated and verified against official GATE question papers and final answer keys.
\par\medskip\noindent
\textbf{Notice on Image-Based Questions:}\\
Questions containing diagrams, circuit schematics, graphs, or image-based figures from the original examination papers have been omitted in this text edition due to non-readable format in historical source archives. All text-based conceptual, analytical, mathematical, algorithmic, and numerical questions are fully included with verified official keys.
}}

\vspace{1.5em}
\begin{center}
\textbf{EXAMINATION STRUCTURE \& PRACTICE GUIDELINES}
\end{center}
\begin{enumerate}[leftmargin=*]
    \item \textbf{Sectional Mark Division:} Each paper comprises General Aptitude (GA) and Computer Science (CS). Questions are clearly grouped into 1-mark and 2-mark blocks at the head of each section.
    \item \textbf{Question Types:} Includes Multiple Choice Questions (MCQ), Multiple Select Questions (MSQ), and Numerical Answer Type (NAT).
    \item \textbf{Official Answer Keys:} For self-evaluation without distraction, complete official answer keys are printed at the conclusion of each respective examination paper set.
\end{enumerate}

\newpage

% ==================== PAGE 3: TABLE OF CONTENTS ====================
\begin{center}
    {\Large\textbf{TABLE OF CONTENTS}}\\[0.4em]
    \rule{0.5\textwidth}{0.6pt}
\end{center}
\vspace{1em}

\noindent\begin{tabularx}{\linewidth}{|X|l|c|r|}
\hline
\textbf{Examination Paper} & \textbf{Session / Shift} & \textbf{Total Questions} & \textbf{Content Type} \\
\hline

GATE 2016 Computer Science & Shift CS & 124 & Questions + Answer Key \\ \hline
\end{tabularx}

\vspace{1.5em}
\noindent{\footnotesize\textit{Note: All questions in this edition have been typeset from verified digital records and cross-referenced with official answer keys released by the organising IIT/IISc institutes.}}

\newpage

% ==================== PAPER: 2016_CS ====================
\begin{center}
{\Large\textbf{GATE 2016 EXAMINATION}}\\[3pt]
{\normalsize\textbf{COMPUTER SCIENCE AND INFORMATION TECHNOLOGY}}\\[2pt]
{\small\textbf{Session / Shift: CS}}\\[5pt]
\rule{\linewidth}{0.8pt}\vspace{-4pt}
{\footnotesize\textbf{Duration: 180 Minutes} \hfill \textbf{Total Marks: 100} \hfill \textbf{Questions: 55}}\\
\vspace{-4pt}\rule{\linewidth}{0.4pt}
\end{center}
\vspace{4pt}
\noindent\fbox{\parbox{\dimexpr\linewidth-2\fboxsep-2\fboxrule}{%
\footnotesize\textbf{Examination Instructions:}\\[1pt]
1. Question Types: Multiple Choice Questions (MCQ), Multiple Select Questions (MSQ), and Numerical Answer Type (NAT).\\
2. Negative Marking Scheme: For 1-mark MCQ, $1/3$ mark is deducted; for 2-mark MCQ, $2/3$ mark is deducted. There is no negative marking for MSQ and NAT questions.\\
3. Mark Division: Questions are grouped by mark allocation within each section as specified below. Full official answer key is provided at the end of this paper.
}}
\vspace{10pt}
\begin{multicols}{2}

\vspace{8pt}
\noindent\textbf{\large\underline{GENERAL APTITUDE}}\\[2pt]
\noindent{\footnotesize\textit{[Section Block: 1-mark questions]}}
\vspace{4pt}

\noindent\textbf{Q.1} \hfill {\footnotesize\textsf{[MCQ]}}
\par\vspace{2pt}
{\bfseries Out of the following four sentences, select the most suitable sentence with respect to grammar and
usage. \par}
\par\vspace{2pt}\noindent
\textbf{(A)}~I will not leave the place until the minister does not meet me.\par\vspace{2pt}\noindent
\textbf{(B)}~I will not leave the place until the minister doesn't meet me.\par\vspace{2pt}\noindent
\textbf{(C)}~I will not leave the place until the minister meet me.\par\vspace{2pt}\noindent
\textbf{(D)}~I will not leave the place until the minister meets me.
\par\vspace{6pt}\hrule height 0.25pt\vspace{6pt}
\noindent\textbf{Q.2} \hfill {\footnotesize\textsf{[MCQ]}}
\par\vspace{2pt}
{\bfseries A rewording of something written or spoken is a \_\_\_\_\_\_\_\_\_\_\_\_\_\_. \par}
\par\vspace{2pt}\noindent
\begin{tabular*}{\columnwidth}{@{\extracolsep{\fill}}p{0.46\columnwidth}p{0.46\columnwidth}@{}}
\textbf{(A)}~paraphrase & \textbf{(B)}~paradox \\
\addlinespace[2pt]
\textbf{(C)}~paradigm & \textbf{(D)}~paraffin
\end{tabular*}
\par\vspace{6pt}\hrule height 0.25pt\vspace{6pt}
\noindent\textbf{Q.3} \hfill {\footnotesize\textsf{[MCQ]}}
\par\vspace{2pt}
{\bfseries Archimedes said, ``Give me a lever long enough and a fulcrum on which to place it, and I will move
the world.''
\par\smallskip
The sentence above is an example of a \_\_\_\_\_\_\_\_\_\_\_ statement. \par}
\par\vspace{2pt}\noindent
\begin{tabular*}{\columnwidth}{@{\extracolsep{\fill}}p{0.46\columnwidth}p{0.46\columnwidth}@{}}
\textbf{(A)}~figurative & \textbf{(B)}~collateral \\
\addlinespace[2pt]
\textbf{(C)}~literal & \textbf{(D)}~figurine
\end{tabular*}
\par\vspace{6pt}\hrule height 0.25pt\vspace{6pt}
\noindent\textbf{Q.4} \hfill {\footnotesize\textsf{[MCQ]}}
\par\vspace{2pt}
{\bfseries If `relftaga' means carefree, `otaga' means careful and `fertaga' means careless, which of the
following could mean `aftercare'? \par}
\par\vspace{2pt}\noindent
\begin{tabular*}{\columnwidth}{@{\extracolsep{\fill}}p{0.46\columnwidth}p{0.46\columnwidth}@{}}
\textbf{(A)}~zentaga & \textbf{(B)}~tagafer \\
\addlinespace[2pt]
\textbf{(C)}~tagazen & \textbf{(D)}~relffer
\end{tabular*}
\par\vspace{6pt}\hrule height 0.25pt\vspace{6pt}
\noindent\textbf{Q.5} \hfill {\footnotesize\textsf{[MCQ]}}
\par\vspace{2pt}
{\bfseries A cube is built using 64 cubic blocks of side one unit. After it is built, one cubic block is removed
from every corner of the cube. The resulting surface area of the body (in square units) after the
removal is \_\_\_\_\_\_\_\_\_\_. \par}
\par\vspace{2pt}\noindent
\begin{tabular*}{\columnwidth}{@{\extracolsep{\fill}}p{0.46\columnwidth}p{0.46\columnwidth}@{}}
\textbf{(A)}~56 & \textbf{(B)}~64 \\
\addlinespace[2pt]
\textbf{(C)}~72 & \textbf{(D)}~96
\end{tabular*}
\par\vspace{6pt}\hrule height 0.25pt\vspace{6pt}

\vspace{6pt}
\noindent{\footnotesize\textit{[Section Block: 2-mark questions]}}
\vspace{3pt}

\noindent\textbf{Q.6} \hfill {\footnotesize\textsf{[MCQ]}}
\par\vspace{2pt}
{\bfseries A shaving set company sells 4 different types of razors, Elegance, Smooth, Soft and Executive.
Elegance sells at Rs. 48, Smooth at Rs. 63, Soft at Rs. 78 and Executive at Rs. 173 per piece. The
table below shows the numbers of each razor sold in each quarter of a year.
Quarter \ Product
Elegance
Smooth
Soft
Executive
Q1
27300
20009
17602
9999
Q2
25222
19392
18445
8942
Q3
28976
22429
19544
10234
Q4
21012
18229
16595
10109
\par\smallskip
Which product contributes the greatest fraction to the revenue of the company in that year? \par}
\par\vspace{2pt}\noindent
\begin{tabular*}{\columnwidth}{@{\extracolsep{\fill}}p{0.46\columnwidth}p{0.46\columnwidth}@{}}
\textbf{(A)}~Elegance & \textbf{(B)}~Executive \\
\addlinespace[2pt]
\textbf{(C)}~Smooth & \textbf{(D)}~Soft
\end{tabular*}
\par\vspace{6pt}\hrule height 0.25pt\vspace{6pt}
\noindent\textbf{Q.7} \hfill {\footnotesize\textsf{[MCQ]}}
\par\vspace{2pt}
{\bfseries Indian currency notes show the denomination indicated in at least seventeen languages. If this is not
an indication of the nation's diversity, nothing else is.
\par\smallskip
Which of the following can be logically inferred from the above sentences? \par}
\par\vspace{2pt}\noindent
\textbf{(A)}~India is a country of exactly seventeen languages.\par\vspace{2pt}\noindent
\textbf{(B)}~Linguistic pluralism is the only indicator of a nation's diversity.\par\vspace{2pt}\noindent
\textbf{(C)}~Indian currency notes have sufficient space for all the Indian languages.\par\vspace{2pt}\noindent
\textbf{(D)}~Linguistic pluralism is strong evidence of India's diversity.
\par\vspace{6pt}\hrule height 0.25pt\vspace{6pt}
\noindent\textbf{Q.8} \hfill {\footnotesize\textsf{[MCQ]}}
\par\vspace{2pt}
{\bfseries Consider the following statements relating to the level of poker play of four players P, Q, R and S.
I.
P always beats Q
II.
R always beats S
III.
S loses to P only sometimes
IV.
R always loses to Q
\par\smallskip
Which of the following can be logically inferred from the above statements?
(i)
P is likely to beat all the three other players
(ii)
S is the absolute worst player in the set \par}
\par\vspace{2pt}\noindent
\begin{tabular*}{\columnwidth}{@{\extracolsep{\fill}}p{0.46\columnwidth}p{0.46\columnwidth}@{}}
\textbf{(A)}~(i) only & \textbf{(B)}~(ii) only \\
\addlinespace[2pt]
\textbf{(C)}~(i) and (ii) & \textbf{(D)}~neither (i) nor (ii)
\end{tabular*}
\par\vspace{6pt}\hrule height 0.25pt\vspace{6pt}
\noindent\textbf{Q.9} \hfill {\footnotesize\textsf{[MCQ]}}
\par\vspace{2pt}
{\bfseries If f($x$$x$) = 2$x$$x$7 + 3$x$$x$−5, which of the following is a factor of f(x)? \par}
\par\vspace{2pt}\noindent
\begin{tabular*}{\columnwidth}{@{\extracolsep{\fill}}p{0.46\columnwidth}p{0.46\columnwidth}@{}}
\textbf{(A)}~(x3+8) & \textbf{(B)}~(x-1) \\
\addlinespace[2pt]
\textbf{(C)}~(2x-5) & \textbf{(D)}~(x+1)
\end{tabular*}
\par\vspace{6pt}\hrule height 0.25pt\vspace{6pt}
\noindent\textbf{Q.10} \hfill {\footnotesize\textsf{[MCQ]}}
\par\vspace{2pt}
{\bfseries In a process, the number of cycles to failure decreases exponentially with an increase in load. At a
load of 80 units, it takes 100 cycles for failure. When the load is halved, it takes 10000 cycles for
failure. The load for which the failure will happen in 5000 cycles is \_\_\_\_\_\_\_\_. \par}
\par\vspace{2pt}\noindent
\begin{tabular*}{\columnwidth}{@{\extracolsep{\fill}}p{0.46\columnwidth}p{0.46\columnwidth}@{}}
\textbf{(A)}~40.00 & \textbf{(B)}~46.02 \\
\addlinespace[2pt]
\textbf{(C)}~60.01 & \textbf{(D)}~92.02
\end{tabular*}
\par\vspace{6pt}\hrule height 0.25pt\vspace{6pt}

\vspace{6pt}
\noindent{\footnotesize\textit{[Section Block: 1-mark questions]}}
\vspace{3pt}

\noindent\textbf{Q.11} \hfill {\footnotesize\textsf{[NAT]}}
\par\vspace{2pt}
{\bfseries Consider the following directed graph:
a
b
c
f
e
d
The number of different topological orderings of the vertices of the graph is
. \par}
\par\vspace{3pt}\noindent\textbf{Answer:} \rule{3.2cm}{0.4pt}
\par\vspace{6pt}\hrule height 0.25pt\vspace{6pt}
\noindent\textbf{Q.12} \hfill {\footnotesize\textsf{[MCQ]}}
\par\vspace{2pt}
{\bfseries Consider the following C program.
void f(int, short);
void main()
\{
int i = 100;
short s = 12;
short *p = \&s;
\_\_\_\_\_\_\_\_\_\_ ;
// call to f()
\}
Which one of the following expressions, when placed in the blank above, will NOT result in
a type checking error? \par}
\par\vspace{2pt}\noindent
\textbf{(A)}~f(s,*s)\par\vspace{2pt}\noindent
\textbf{(B)}~i = f(i,s)\par\vspace{2pt}\noindent
\textbf{(C)}~f(i,*s)\par\vspace{2pt}\noindent
\textbf{(D)}~f(i,*p) \newline CS(Set A)
\par\vspace{6pt}\hrule height 0.25pt\vspace{6pt}
\noindent\textbf{Q.13} \hfill {\footnotesize\textsf{[MCQ]}}
\par\vspace{2pt}
{\bfseries The worst case running times of Insertion sort, Merge sort and Quick sort, respectively, are: \par}
\par\vspace{2pt}\noindent
\textbf{(A)}~Θ(nlogn), Θ(nlogn), and Θ(n2)\par\vspace{2pt}\noindent
\textbf{(B)}~Θ(n2), Θ(n2), and Θ(nlogn)\par\vspace{2pt}\noindent
\textbf{(C)}~Θ(n2), Θ(nlogn), and Θ(nlogn)\par\vspace{2pt}\noindent
\textbf{(D)}~Θ(n2), Θ(nlogn), and Θ(n2)
\par\vspace{6pt}\hrule height 0.25pt\vspace{6pt}
\noindent\textbf{Q.14} \hfill {\footnotesize\textsf{[MCQ]}}
\par\vspace{2pt}
{\bfseries Let G be a weighted connected undirected graph with distinct positive edge weights. If every
edge weight is increased by the same value, then which of the following statements is/are
TRUE?
P: Minimum spanning tree of G does not change
Q: Shortest path between any pair of vertices does not change \par}
\par\vspace{2pt}\noindent
\textbf{(A)}~P only\par\vspace{2pt}\noindent
\textbf{(B)}~Q only\par\vspace{2pt}\noindent
\textbf{(C)}~Neither P nor Q\par\vspace{2pt}\noindent
\textbf{(D)}~Both P and Q \newline CS(Set A)
\par\vspace{6pt}\hrule height 0.25pt\vspace{6pt}
\noindent\textbf{Q.15} \hfill {\footnotesize\textsf{[NAT]}}
\par\vspace{2pt}
{\bfseries Consider the following C program.
\#include<stdio.h>
void mystery(int *ptra, int *ptrb) \{
int *temp;
temp = ptrb;
ptrb = ptra;
ptra = temp;
\}
int main() \{
int a=2016, b=0, c=4, d=42;
mystery(\&a, \&b);
if (a < c)
mystery(\&c, \&a);
mystery(\&a, \&d);
printf("\%d\n", a);
\}
The output of the program is
. \par}
\par\vspace{3pt}\noindent\textbf{Answer:} \rule{3.2cm}{0.4pt}
\par\vspace{6pt}\hrule height 0.25pt\vspace{6pt}
\noindent\textbf{Q.16} \hfill {\footnotesize\textsf{[MCQ]}}
\par\vspace{2pt}
{\bfseries Which of the following languages is generated by the given grammar?
S
−$\to$
aS | bS | $\epsilon$ \par}
\par\vspace{2pt}\noindent
\textbf{(A)}~\{anbm | n,m $\ge$0\}\par\vspace{2pt}\noindent
\textbf{(B)}~\{w $\in$\{a,b\}∗| w has equal number of a's and b's \}\par\vspace{2pt}\noindent
\textbf{(C)}~\{an | n $\ge$0\}$\cup$\{bn | n $\ge$0\}$\cup$\{anbn | n $\ge$0\}\par\vspace{2pt}\noindent
\textbf{(D)}~\{a,b\}∗ \newline CS(Set A)
\par\vspace{6pt}\hrule height 0.25pt\vspace{6pt}
\noindent\textbf{Q.17} \hfill {\footnotesize\textsf{[MCQ]}}
\par\vspace{2pt}
{\bfseries Which of the following decision problems are undecidable?
I. Given NFAs N1 and N2, is L(N1)$\cap$L(N2) = Φ?
II. Given a CFG G = (N,$\Sigma$,P,S) and a string x $\in$$\Sigma$∗, does x $\in$L(G)?
III. Given CFGs G1 and G2, is L(G1) = L(G2)?
IV. Given a TM M, is L(M) = Φ? \par}
\par\vspace{2pt}\noindent
\begin{tabular*}{\columnwidth}{@{\extracolsep{\fill}}p{0.46\columnwidth}p{0.46\columnwidth}@{}}
\textbf{(A)}~I and IV only & \textbf{(B)}~II and III only \\
\addlinespace[2pt]
\textbf{(C)}~III and IV only & \textbf{(D)}~II and IV only
\end{tabular*}
\par\vspace{6pt}\hrule height 0.25pt\vspace{6pt}
\noindent\textbf{Q.18} \hfill {\footnotesize\textsf{[MCQ]}}
\par\vspace{2pt}
{\bfseries Which one of the following regular expressions represents the language: the set of all binary
strings having two consecutive 0s and two consecutive 1s? \par}
\par\vspace{2pt}\noindent
\textbf{(A)}~(0+1)∗0011(0+1)∗+(0+1)∗1100(0+1)∗\par\vspace{2pt}\noindent
\textbf{(B)}~(0+1)∗(00(0+1)∗11+11(0+1)∗00)(0+1)∗\par\vspace{2pt}\noindent
\textbf{(C)}~(0+1)∗00(0+1)∗+(0+1)∗11(0+1)∗\par\vspace{2pt}\noindent
\textbf{(D)}~00(0+1)∗11+11(0+1)∗00
\par\vspace{6pt}\hrule height 0.25pt\vspace{6pt}
\noindent\textbf{Q.19} \hfill {\footnotesize\textsf{[NAT]}}
\par\vspace{2pt}
{\bfseries Consider the following code segment.
x = u - t;
y = x * v;
x = y + w;
y = t - z;
y = x * y;
The minimum number of total variables required to convert the above code segment to static
single assignment form is
.
CS(Set A) \par}
\par\vspace{3pt}\noindent\textbf{Answer:} \rule{3.2cm}{0.4pt}
\par\vspace{6pt}\hrule height 0.25pt\vspace{6pt}
\noindent\textbf{Q.20} \hfill {\footnotesize\textsf{[MCQ]}}
\par\vspace{2pt}
{\bfseries Consider an arbitrary set of CPU-bound processes with unequal CPU burst lengths
submitted at the same time to a computer system.
Which one of the following process
scheduling algorithms would minimize the average waiting time in the ready queue? \par}
\par\vspace{2pt}\noindent
\textbf{(A)}~Shortest remaining time first\par\vspace{2pt}\noindent
\textbf{(B)}~Round-robin with time quantum less than the shortest CPU burst\par\vspace{2pt}\noindent
\textbf{(C)}~Uniform random\par\vspace{2pt}\noindent
\textbf{(D)}~Highest priority first with priority proportional to CPU burst length
\par\vspace{6pt}\hrule height 0.25pt\vspace{6pt}
\noindent\textbf{Q.21} \hfill {\footnotesize\textsf{[MCQ]}}
\par\vspace{2pt}
{\bfseries Which of the following is NOT a superkey in a relational schema with attributes
V, W, X, Y, Z and primary key V Y? \par}
\par\vspace{2pt}\noindent
\begin{tabular*}{\columnwidth}{@{\extracolsep{\fill}}p{0.46\columnwidth}p{0.46\columnwidth}@{}}
\textbf{(A)}~VXYZ & \textbf{(B)}~VWXZ \\
\addlinespace[2pt]
\textbf{(C)}~VWXY & \textbf{(D)}~VWXYZ
\end{tabular*}
\par\vspace{6pt}\hrule height 0.25pt\vspace{6pt}
\noindent\textbf{Q.22} \hfill {\footnotesize\textsf{[MCQ]}}
\par\vspace{2pt}
{\bfseries Which one of the following is NOT a part of the ACID properties of database transactions? \par}
\par\vspace{2pt}\noindent
\textbf{(A)}~Atomicity\par\vspace{2pt}\noindent
\textbf{(B)}~Consistency\par\vspace{2pt}\noindent
\textbf{(C)}~Isolation\par\vspace{2pt}\noindent
\textbf{(D)}~Deadlock-freedom \newline CS(Set A)
\par\vspace{6pt}\hrule height 0.25pt\vspace{6pt}
\noindent\textbf{Q.23} \hfill {\footnotesize\textsf{[MCQ]}}
\par\vspace{2pt}
{\bfseries A database of research articles in a journal uses the following schema.
(VOLUME, NUMBER, STARTPAGE, ENDPAGE, TITLE, YEAR, PRICE)
The primary key is (VOLUME, NUMBER, STARTPAGE, ENDPAGE) and the following
functional dependencies exist in the schema.
(VOLUME, NUMBER, STARTPAGE, ENDPAGE)
$\to$
TITLE
(VOLUME, NUMBER)
$\to$
YEAR
(VOLUME, NUMBER, STARTPAGE, ENDPAGE)
$\to$
PRICE
The database is redesigned to use the following schemas.
(VOLUME, NUMBER, STARTPAGE, ENDPAGE, TITLE, PRICE)
(VOLUME, NUMBER, YEAR)
Which is the weakest normal form that the new database satisfies, but the old one does not? \par}
\par\vspace{2pt}\noindent
\begin{tabular*}{\columnwidth}{@{\extracolsep{\fill}}p{0.46\columnwidth}p{0.46\columnwidth}@{}}
\textbf{(A)}~1NF & \textbf{(B)}~2NF \\
\addlinespace[2pt]
\textbf{(C)}~3NF & \textbf{(D)}~BCNF
\end{tabular*}
\par\vspace{6pt}\hrule height 0.25pt\vspace{6pt}
\noindent\textbf{Q.24} \hfill {\footnotesize\textsf{[MCQ]}}
\par\vspace{2pt}
{\bfseries Which one of the following protocols is NOT used to resolve one form of address to another
one? \par}
\par\vspace{2pt}\noindent
\begin{tabular*}{\columnwidth}{@{\extracolsep{\fill}}p{0.46\columnwidth}p{0.46\columnwidth}@{}}
\textbf{(A)}~DNS & \textbf{(B)}~ARP \\
\addlinespace[2pt]
\textbf{(C)}~DHCP & \textbf{(D)}~RARP
\end{tabular*}
\par\vspace{6pt}\hrule height 0.25pt\vspace{6pt}
\noindent\textbf{Q.25} \hfill {\footnotesize\textsf{[MCQ]}}
\par\vspace{2pt}
{\bfseries Which of the following is/are example(s) of stateful application layer protocols?
(i) HTTP
(ii) FTP
(iii) TCP
(iv) POP3 \par}
\par\vspace{2pt}\noindent
\textbf{(A)}~(i) and (ii) only\par\vspace{2pt}\noindent
\textbf{(B)}~(ii) and (iii) only\par\vspace{2pt}\noindent
\textbf{(C)}~(ii) and (iv) only\par\vspace{2pt}\noindent
\textbf{(D)}~(iv) only \newline CS(Set A)
\par\vspace{6pt}\hrule height 0.25pt\vspace{6pt}

\vspace{6pt}
\noindent{\footnotesize\textit{[Section Block: 2-mark questions]}}
\vspace{3pt}

\noindent\textbf{Q.26} \hfill {\footnotesize\textsf{[NAT]}}
\par\vspace{2pt}
{\bfseries The coefficient of x12 in (x3 +x4 +x5 +x6 +···)3 is
. \par}
\par\vspace{3pt}\noindent\textbf{Answer:} \rule{3.2cm}{0.4pt}
\par\vspace{6pt}\hrule height 0.25pt\vspace{6pt}
\noindent\textbf{Q.27} \hfill {\footnotesize\textsf{[NAT]}}
\par\vspace{2pt}
{\bfseries Consider the recurrence relation a1 = 8, an = 6n2 +2n+an−1. Let a99 = K $\times$104. The value
of K is
. \par}
\par\vspace{3pt}\noindent\textbf{Answer:} \rule{3.2cm}{0.4pt}
\par\vspace{6pt}\hrule height 0.25pt\vspace{6pt}
\noindent\textbf{Q.28} \hfill {\footnotesize\textsf{[NAT]}}
\par\vspace{2pt}
{\bfseries A function f : N+ $\to$N+, defined on the set of positive integers N+, satisfies the following
properties:
f(n) = f(n/2)
if nis even
f(n) = f(n+5)
if nis odd
Let R = \{i|∃j : f(j) = i\} be the set of distinct values that f takes. The maximum possible size
of R is
. \par}
\par\vspace{3pt}\noindent\textbf{Answer:} \rule{3.2cm}{0.4pt}
\par\vspace{6pt}\hrule height 0.25pt\vspace{6pt}
\noindent\textbf{Q.29} \hfill {\footnotesize\textsf{[NAT]}}
\par\vspace{2pt}
{\bfseries Consider the following experiment.
Step 1. Flip a fair coin twice.
Step 2. If the outcomes are (TAILS, HEADS) then output Y and stop.
Step 3. If the outcomes are either (HEADS, HEADS) or (HEADS, TAILS), then output N and
stop.
Step 4. If the outcomes are (TAILS, TAILS), then go to Step 1.
The probability that the output of the experiment is Y is (up to two decimal places)
.
CS(Set A) \par}
\par\vspace{3pt}\noindent\textbf{Answer:} \rule{3.2cm}{0.4pt}
\par\vspace{6pt}\hrule height 0.25pt\vspace{6pt}
\noindent\textbf{Q.30} \hfill {\footnotesize\textsf{[NAT]}}
\par\vspace{2pt}
{\bfseries Suppose the functions F and G can be computed in 5 and 3 nanoseconds by functional units
UF and UG, respectively. Given two instances of UF and two instances of UG, it is required to
implement the computation F(G(Xi)) for 1 $\le$i $\le$10. Ignoring all other delays, the minimum
time required to complete this computation is
nanoseconds. \par}
\par\vspace{3pt}\noindent\textbf{Answer:} \rule{3.2cm}{0.4pt}
\par\vspace{6pt}\hrule height 0.25pt\vspace{6pt}
\noindent\textbf{Q.31} \hfill {\footnotesize\textsf{[NAT]}}
\par\vspace{2pt}
{\bfseries The size of the data count register of a DMA controller is 16 bits. The processor needs to
transfer a file of 29,154 kilobytes from disk to main memory. The memory is byte addressable.
The minimum number of times the DMA controller needs to get the control of the system bus
from the processor to transfer the file from the disk to main memory is
. \par}
\par\vspace{3pt}\noindent\textbf{Answer:} \rule{3.2cm}{0.4pt}
\par\vspace{6pt}\hrule height 0.25pt\vspace{6pt}
\noindent\textbf{Q.32} \hfill {\footnotesize\textsf{[NAT]}}
\par\vspace{2pt}
{\bfseries The stage delays in a 4-stage pipeline are 800, 500, 400 and 300 picoseconds. The first
stage (with delay 800 picoseconds) is replaced with a functionally equivalent design involving
two stages with respective delays 600 and 350 picoseconds. The throughput increase of the
pipeline is
percent. \par}
\par\vspace{3pt}\noindent\textbf{Answer:} \rule{3.2cm}{0.4pt}
\par\vspace{6pt}\hrule height 0.25pt\vspace{6pt}
\noindent\textbf{Q.33} \hfill {\footnotesize\textsf{[MCQ]}}
\par\vspace{2pt}
{\bfseries Consider a carry lookahead adder for adding two n-bit integers, built using gates of fan-in at
most two. The time to perform addition using this adder is \par}
\par\vspace{2pt}\noindent
\textbf{(A)}~Θ(1)\par\vspace{2pt}\noindent
\textbf{(B)}~Θ(log(n))\par\vspace{2pt}\noindent
\textbf{(C)}~Θ($\sqrt{\vphantom{x}}$n)\par\vspace{2pt}\noindent
\textbf{(D)}~Θ(n) \newline CS(Set A)
\par\vspace{6pt}\hrule height 0.25pt\vspace{6pt}
\noindent\textbf{Q.34} \hfill {\footnotesize\textsf{[MCQ]}}
\par\vspace{2pt}
{\bfseries The following function computes the maximum value contained in an integer array p[] of size
n (n $\ge$ 1).
int max(int *p, int n) \{
int a=0, b=n-1;
while (\_\_\_\_\_\_\_\_\_\_) \{
if (p[a] $\le$ p[b]) \{ a = a+1; \}
else
\{ b = b-1; \}
\}
return p[a];
\}
The missing loop condition is \par}
\par\vspace{2pt}\noindent
\begin{tabular*}{\columnwidth}{@{\extracolsep{\fill}}p{0.46\columnwidth}p{0.46\columnwidth}@{}}
\textbf{(A)}~a $\ne$ n & \textbf{(B)}~b $\ne$ 0 \\
\addlinespace[2pt]
\textbf{(C)}~b > (a + 1) & \textbf{(D)}~b $\ne$ a
\end{tabular*}
\par\vspace{6pt}\hrule height 0.25pt\vspace{6pt}
\noindent\textbf{Q.35} \hfill {\footnotesize\textsf{[MCQ]}}
\par\vspace{2pt}
{\bfseries What will be the output of the following C program?
void count(int n)\{
static int d=1;
printf("\%d ", n);
printf("\%d ", d);
d++;
if(n>1) count(n-1);
printf("\%d ", d);
\}
void main()\{
count(3);
\} \par}
\par\vspace{2pt}\noindent
\textbf{(A)}~3 1 2 2 1 3 4 4 4\par\vspace{2pt}\noindent
\textbf{(B)}~3 1 2 1 1 1 2 2 2\par\vspace{2pt}\noindent
\textbf{(C)}~3 1 2 2 1 3 4\par\vspace{2pt}\noindent
\textbf{(D)}~3 1 2 1 1 1 2 \newline CS(Set A)
\par\vspace{6pt}\hrule height 0.25pt\vspace{6pt}
\noindent\textbf{Q.36} \hfill {\footnotesize\textsf{[MCQ]}}
\par\vspace{2pt}
{\bfseries What will be the output of the following pseudo-code when parameters are passed by reference
and dynamic scoping is assumed?
a=3;
void n(x) \{x = x * a; print(x);\}
void m(y) \{a = 1; a = y - a; n(a); print(a);\}
void main() \{m(a);\} \par}
\par\vspace{2pt}\noindent
\begin{tabular*}{\columnwidth}{@{\extracolsep{\fill}}p{0.46\columnwidth}p{0.46\columnwidth}@{}}
\textbf{(A)}~6, 2 & \textbf{(B)}~6, 6 \\
\addlinespace[2pt]
\textbf{(C)}~4, 2 & \textbf{(D)}~4, 4
\end{tabular*}
\par\vspace{6pt}\hrule height 0.25pt\vspace{6pt}
\noindent\textbf{Q.37} \hfill {\footnotesize\textsf{[MCQ]}}
\par\vspace{2pt}
{\bfseries An operator delete(i) for a binary heap data structure is to be designed to delete the item in
the i-th node. Assume that the heap is implemented in an array and i refers to the i-th index
of the array. If the heap tree has depth d (number of edges on the path from the root to the
farthest leaf), then what is the time complexity to re-fix the heap efficiently after the removal
of the element? \par}
\par\vspace{2pt}\noindent
\begin{tabular*}{\columnwidth}{@{\extracolsep{\fill}}p{0.46\columnwidth}p{0.46\columnwidth}@{}}
\textbf{(A)}~O(1) & \textbf{(B)}~O(d) but not O(1) \\
\addlinespace[2pt]
\textbf{(C)}~O(2d) but not O(d) & \textbf{(D)}~O(d 2d) but not O(2d)
\end{tabular*}
\par\vspace{6pt}\hrule height 0.25pt\vspace{6pt}
\noindent\textbf{Q.38} \hfill {\footnotesize\textsf{[NAT]}}
\par\vspace{2pt}
{\bfseries Consider the weighted undirected graph with 4 vertices, where the weight of edge \{i, j\} is
given by the entry Wi j in the matrix W.
W =


0
2
8
5
2
0
5
8
8
5
0
x
5
8
x
0


The largest possible integer value of x, for which at least one shortest path between some pair
of vertices will contain the edge with weight x is
. \par}
\par\vspace{3pt}\noindent\textbf{Answer:} \rule{3.2cm}{0.4pt}
\par\vspace{6pt}\hrule height 0.25pt\vspace{6pt}
\noindent\textbf{Q.39} \hfill {\footnotesize\textsf{[NAT]}}
\par\vspace{2pt}
{\bfseries Let G be a complete undirected graph on 4 vertices, having 6 edges with weights being 1, 2,
3, 4, 5, and 6. The maximum possible weight that a minimum weight spanning tree of G can
have is
.
CS(Set A) \par}
\par\vspace{3pt}\noindent\textbf{Answer:} \rule{3.2cm}{0.4pt}
\par\vspace{6pt}\hrule height 0.25pt\vspace{6pt}
\noindent\textbf{Q.40} \hfill {\footnotesize\textsf{[MCQ]}}
\par\vspace{2pt}
{\bfseries G = (V,E) is an undirected simple graph in which each edge has a distinct weight, and e is a
particular edge of G. Which of the following statements about the minimum spanning trees
(MSTs) of G is/are TRUE?
I. If e is the lightest edge of some cycle in G, then every MST of G includes e
II. If e is the heaviest edge of some cycle in G, then every MST of G excludes e \par}
\par\vspace{2pt}\noindent
\begin{tabular*}{\columnwidth}{@{\extracolsep{\fill}}p{0.46\columnwidth}p{0.46\columnwidth}@{}}
\textbf{(A)}~I only & \textbf{(B)}~II only \\
\addlinespace[2pt]
\textbf{(C)}~both I and II & \textbf{(D)}~neither I nor II
\end{tabular*}
\par\vspace{6pt}\hrule height 0.25pt\vspace{6pt}
\noindent\textbf{Q.41} \hfill {\footnotesize\textsf{[NAT]}}
\par\vspace{2pt}
{\bfseries Let Q denote a queue containing sixteen numbers and S be an empty stack. Head(Q) returns
the element at the head of the queue Q without removing it from Q. Similarly Top(S) returns
the element at the top of S without removing it from S. Consider the algorithm given below.
while Q is not Empty do
if S is Empty OR Top(S) $\le$Head(Q) then
x := Dequeue(Q);
Push(S,x);
else
x := Pop(S);
Enqueue(Q,x);
end
end
The maximum possible number of iterations of the while loop in the algorithm is
. \par}
\par\vspace{3pt}\noindent\textbf{Answer:} \rule{3.2cm}{0.4pt}
\par\vspace{6pt}\hrule height 0.25pt\vspace{6pt}
\noindent\textbf{Q.42} \hfill {\footnotesize\textsf{[MCQ]}}
\par\vspace{2pt}
{\bfseries Consider the following context-free grammars:
G1: S $\to$aS|B, B $\to$b|bB
G2: S $\to$aA|bB, A $\to$aA|B|$\epsilon$, B $\to$bB|$\epsilon$
Which one of the following pairs of languages is generated by G1 and G2, respectively? \par}
\par\vspace{2pt}\noindent
\textbf{(A)}~\{ambn|m > 0 or n > 0\} and \{ambn|m > 0 and n > 0\}\par\vspace{2pt}\noindent
\textbf{(B)}~\{ambn|m > 0 and n > 0\} and \{ambn|m > 0 or n $\ge$0\}\par\vspace{2pt}\noindent
\textbf{(C)}~\{ambn|m $\ge$0 or n > 0\} and \{ambn|m > 0 and n > 0\}\par\vspace{2pt}\noindent
\textbf{(D)}~\{ambn|m $\ge$0 and n > 0\} and \{ambn|m > 0 or n > 0\} \newline CS(Set A)
\par\vspace{6pt}\hrule height 0.25pt\vspace{6pt}
\noindent\textbf{Q.43} \hfill {\footnotesize\textsf{[MCQ]}}
\par\vspace{2pt}
{\bfseries Consider the following languages:
L1 = \{anbmcn+m : m,n $\ge$1\}
L2 = \{anbnc2n : n $\ge$1\}
Which one of the following is TRUE? \par}
\par\vspace{2pt}\noindent
\textbf{(A)}~Both L1 and L2 are context-free.\par\vspace{2pt}\noindent
\textbf{(B)}~L1 is context-free while L2 is not context-free.\par\vspace{2pt}\noindent
\textbf{(C)}~L2 is context-free while L1 is not context-free.\par\vspace{2pt}\noindent
\textbf{(D)}~Neither L1 nor L2 is context-free. \newline CS(Set B)
\par\vspace{6pt}\hrule height 0.25pt\vspace{6pt}
\noindent\textbf{Q.44} \hfill {\footnotesize\textsf{[MCQ]}}
\par\vspace{2pt}
{\bfseries Let X be a recursive language and Y be a recursively enumerable but not recursive language.
Let W and Z be two languages such that Y reduces to W, and Z reduces to X (reduction means
the standard many-one reduction). Which one of the following statements is TRUE? \par}
\par\vspace{2pt}\noindent
\textbf{(A)}~W can be recursively enumerable and Z is recursive.\par\vspace{2pt}\noindent
\textbf{(B)}~W can be recursive and Z is recursively enumerable.\par\vspace{2pt}\noindent
\textbf{(C)}~W is not recursively enumerable and Z is recursive.\par\vspace{2pt}\noindent
\textbf{(D)}~W is not recursively enumerable and Z is not recursive.
\par\vspace{6pt}\hrule height 0.25pt\vspace{6pt}
\noindent\textbf{Q.45} \hfill {\footnotesize\textsf{[NAT]}}
\par\vspace{2pt}
{\bfseries The attributes of three arithmetic operators in some programming language are given below.
Operator
Precedence
Associativity
Arity
+
High
Left
Binary
−
Medium
Right
Binary
∗
Low
Left
Binary
The value of the expression 2−5+1−7∗3 in this language is
.
CS(Set A) \par}
\par\vspace{3pt}\noindent\textbf{Answer:} \rule{3.2cm}{0.4pt}
\par\vspace{6pt}\hrule height 0.25pt\vspace{6pt}
\noindent\textbf{Q.46} \hfill {\footnotesize\textsf{[MCQ]}}
\par\vspace{2pt}
{\bfseries Consider the following Syntax Directed Translation Scheme (SDTS), with non-terminals
\{S, A\} and terminals \{a, b\}.
S
−$\to$
aA
\{ print 1 \}
S
−$\to$
a
\{ print 2 \}
A
−$\to$
Sb
\{ print 3 \}
Using the above SDTS, the output printed by a bottom-up parser, for the input aab is: \par}
\par\vspace{2pt}\noindent
\begin{tabular*}{\columnwidth}{@{\extracolsep{\fill}}p{0.46\columnwidth}p{0.46\columnwidth}@{}}
\textbf{(A)}~1 3 2 & \textbf{(B)}~2 2 3 \\
\addlinespace[2pt]
\textbf{(C)}~2 3 1 & \textbf{(D)}~syntax error
\end{tabular*}
\par\vspace{6pt}\hrule height 0.25pt\vspace{6pt}
\noindent\textbf{Q.47} \hfill {\footnotesize\textsf{[NAT]}}
\par\vspace{2pt}
{\bfseries Consider a computer system with 40-bit virtual addressing and page size of sixteen
kilobytes. If the computer system has a one-level page table per process and each page table
entry requires 48 bits, then the size of the per-process page table is
megabytes. \par}
\par\vspace{3pt}\noindent\textbf{Answer:} \rule{3.2cm}{0.4pt}
\par\vspace{6pt}\hrule height 0.25pt\vspace{6pt}
\noindent\textbf{Q.48} \hfill {\footnotesize\textsf{[NAT]}}
\par\vspace{2pt}
{\bfseries Consider a disk queue with requests for I/O to blocks on cylinders 47, 38, 121, 191, 87, 11,
92, 10. The C-LOOK scheduling algorithm is used. The head is initially at cylinder number
63, moving towards larger cylinder numbers on its servicing pass. The cylinders are numbered
from 0 to 199. The total head movement (in number of cylinders) incurred while servicing
these requests is
. \par}
\par\vspace{3pt}\noindent\textbf{Answer:} \rule{3.2cm}{0.4pt}
\par\vspace{6pt}\hrule height 0.25pt\vspace{6pt}
\noindent\textbf{Q.49} \hfill {\footnotesize\textsf{[NAT]}}
\par\vspace{2pt}
{\bfseries Consider a computer system with ten physical page frames. The system is provided with
an access sequence (a1,a2,...,a20,a1,a2,...,a20), where each ai is a distinct virtual page
number.
The difference in the number of page faults between the last-in-first-out page
replacement policy and the optimal page replacement policy is
.
CS(Set A) \par}
\par\vspace{3pt}\noindent\textbf{Answer:} \rule{3.2cm}{0.4pt}
\par\vspace{6pt}\hrule height 0.25pt\vspace{6pt}
\noindent\textbf{Q.50} \hfill {\footnotesize\textsf{[MCQ]}}
\par\vspace{2pt}
{\bfseries Consider the following proposed solution for the critical section problem.
There are n
processes: P0 ...Pn−1. In the code, function pmax returns an integer not smaller than any
of its arguments. For all i, t[i] is initialized to zero.
Code for Pi:
do \{
c[i]=1; t[i] = pmax(t[0],...,t[n-1])+1; c[i]=0;
for every j̸ = i in \{0,...,n-1\} \{
while (c[j]);
while (t[j] $\ne$ 0 \&\& t[j]$\le$t[i]);
\}
Critical Section;
t[i]=0;
Remainder Section;
\} while (true);
Which one of the following is TRUE about the above solution? \par}
\par\vspace{2pt}\noindent
\textbf{(A)}~At most one process can be in the critical section at any time\par\vspace{2pt}\noindent
\textbf{(B)}~The bounded wait condition is satisfied\par\vspace{2pt}\noindent
\textbf{(C)}~The progress condition is satisfied\par\vspace{2pt}\noindent
\textbf{(D)}~It cannot cause a deadlock
\par\vspace{6pt}\hrule height 0.25pt\vspace{6pt}
\noindent\textbf{Q.51} \hfill {\footnotesize\textsf{[MCQ]}}
\par\vspace{2pt}
{\bfseries Consider the following two phase locking protocol. Suppose a transaction T accesses (for
read or write operations), a certain set of objects \{O1,...,Ok\}. This is done in the following
manner:
Step 1. T acquires exclusive locks to O1, ..., Ok in increasing order of their addresses.
Step 2. The required operations are performed.
Step 3. All locks are released.
This protocol will \par}
\par\vspace{2pt}\noindent
\textbf{(A)}~guarantee serializability and deadlock-freedom\par\vspace{2pt}\noindent
\textbf{(B)}~guarantee neither serializability nor deadlock-freedom\par\vspace{2pt}\noindent
\textbf{(C)}~guarantee serializability but not deadlock-freedom\par\vspace{2pt}\noindent
\textbf{(D)}~guarantee deadlock-freedom but not serializability \newline CS(Set A)
\par\vspace{6pt}\hrule height 0.25pt\vspace{6pt}
\noindent\textbf{Q.52} \hfill {\footnotesize\textsf{[MCQ]}}
\par\vspace{2pt}
{\bfseries Consider that B wants to send a message m that is digitally signed to A. Let the pair of private
and public keys for A and B be denoted by K−
x and K+
x for x = A,B, respectively. Let Kx(m)
represent the operation of encrypting m with a key Kx and H(m) represent the message digest.
Which one of the following indicates the CORRECT way of sending the message m along
with the digital signature to A? \par}
\par\vspace{2pt}\noindent
\textbf{(A)}~\{m,K+ \newline B (H(m))\}\par\vspace{2pt}\noindent
\textbf{(B)}~\{m,K− \newline B (H(m))\}\par\vspace{2pt}\noindent
\textbf{(C)}~\{m,K− \newline A (H(m))\}\par\vspace{2pt}\noindent
\textbf{(D)}~\{m,K+ \newline A (m)\}
\par\vspace{6pt}\hrule height 0.25pt\vspace{6pt}
\noindent\textbf{Q.53} \hfill {\footnotesize\textsf{[NAT]}}
\par\vspace{2pt}
{\bfseries An IP datagram of size 1000 bytes arrives at a router. The router has to forward this packet on
a link whose MTU (maximum transmission unit) is 100 bytes. Assume that the size of the IP
header is 20 bytes.
The number of fragments that the IP datagram will be divided into for transmission is
. \par}
\par\vspace{3pt}\noindent\textbf{Answer:} \rule{3.2cm}{0.4pt}
\par\vspace{6pt}\hrule height 0.25pt\vspace{6pt}
\noindent\textbf{Q.54} \hfill {\footnotesize\textsf{[NAT]}}
\par\vspace{2pt}
{\bfseries For a host machine that uses the token bucket algorithm for congestion control, the token
bucket has a capacity of 1 megabyte and the maximum output rate is 20 megabytes per second.
Tokens arrive at a rate to sustain output at a rate of 10 megabytes per second. The token bucket
is currently full and the machine needs to send 12 megabytes of data. The minimum time
required to transmit the data is
seconds. \par}
\par\vspace{3pt}\noindent\textbf{Answer:} \rule{3.2cm}{0.4pt}
\par\vspace{6pt}\hrule height 0.25pt\vspace{6pt}
\noindent\textbf{Q.55} \hfill {\footnotesize\textsf{[NAT]}}
\par\vspace{2pt}
{\bfseries A sender uses the Stop-and-Wait ARQ protocol for reliable transmission of frames. Frames
are of size 1000 bytes and the transmission rate at the sender is 80 Kbps (1Kbps = 1000
bits/second). Size of an acknowledgement is 100 bytes and the transmission rate at the receiver
is 8 Kbps. The one-way propagation delay is 100 milliseconds.
Assuming no frame is lost, the sender throughput is
bytes/second.
CS(Set A)
\par\smallskip
Q. No Type
Section
Key
Marks
1
MCQ
GA
D
1
2
MCQ
GA
A
1
3
MCQ
GA
A
1
4
MCQ
GA
C
1
5
MCQ
GA
D
1
6
MCQ
GA
B
2
7
MCQ
GA
D
2
8
MCQ
GA
D
2
9
MCQ
GA
B
2
10
MCQ
GA
B
2
1
NAT
CS-1
11.0 : 11.0
1
2
MCQ
CS-1
B
1
3
NAT
CS-1
1.0 : 1.0
1
4
NAT
CS-1
0.5 : 0.5
1
5
NAT
CS-1
15.0 : 15.0
1
6
MCQ
CS-1
A
1
7
NAT
CS-1
-11.0 : -11.0
1
8
NAT
CS-1
3.0 : 4.0
1
9
NAT
CS-1
31.0 : 31.0
1
10
MCQ
CS-1
A
1
11
NAT
CS-1
6.0 : 6.0
1
12
MCQ
CS-1
D
1
13
MCQ
CS-1
D
1
14
MCQ
CS-1
A
1
15
NAT
CS-1
2016.0 : 2016.0
1
16
MCQ
CS-1
D
1
17
MCQ
CS-1
C
1
18
MCQ
CS-1
B
1
19
NAT
CS-1
10.0 : 10.0
1
20
MCQ
CS-1
A
1
21
MCQ
CS-1
B
1
22
MCQ
CS-1
D
1
23
MCQ
CS-1
B
1
24
MCQ
CS-1
C
1
25
MCQ
CS-1
C
1
26
NAT
CS-1
10.0 : 10.0
2
27
NAT
CS-1
197.9 : 198.1
2
28
NAT
CS-1
2.0 : 2.0
2
29
NAT
CS-1
0.33 : 0.34
2
30
MCQ
CS-1
D
2
31
NAT
CS-1
456.0 : 456.0
2
32
NAT
CS-1
33.0 : 34.0
2
33
MCQ
CS-1
B
2
34
MCQ
CS-1
D
2
35
MCQ
CS-1
A
2
36
MCQ
CS-1
D
2
37
MCQ
CS-1
B
2
38
NAT
CS-1
12.0 : 12.0
2
39
NAT
CS-1
7.0 : 7.0
2
40
MCQ
CS-1
B
2
41
NAT
CS-1
256.0 : 256.0
2
42
MCQ
CS-1
D
2
43
MCQ
CS-1
D
2
44
MCQ
CS-1
C
2
45
NAT
CS-1
9.0 : 9.0
2
46
MCQ
CS-1
C
2
47
NAT
CS-1
384.0 : 384.0
2
48
NAT
CS-1
346.0 : 346.0
2
49
NAT
CS-1
1.0 : 1.0
2
50
MCQ
CS-1
A
2
51
MCQ
CS-1
A
2
52
MCQ
CS-1
B
2
53
NAT
CS-1
13.0 : 13.0
2
54
NAT
CS-1
1.1 : 1.1
2
55
NAT
CS-1
2500.0 : 2500.0
2 \par}
\par\vspace{3pt}\noindent\textbf{Answer:} \rule{3.2cm}{0.4pt}
\par\vspace{6pt}\hrule height 0.25pt\vspace{6pt}
\end{multicols}
\vspace{8pt}
\begin{center}
\rule{0.35\linewidth}{0.5pt}\\[2pt]
\textbf{\small --- END OF QUESTION PAPER ---}\\[-8pt]
\rule{0.35\linewidth}{0.5pt}
\end{center}
\vspace{10pt}
\subsection*{OFFICIAL ANSWER KEY --- GATE 2016 (CS)}
\vspace{-4pt}
\noindent{\footnotesize\textit{Official final answer keys published by the GATE Organising Institute. For NAT questions, allowable numerical ranges are displayed where applicable. MTA = Marks To All.}}
\vspace{4pt}
\begin{center}
\small
\resizebox{\linewidth}{!}{%
\begin{tabular}{|c|c|l||c|c|l||c|c|l||c|c|l|}
\hline
\textbf{Q} & \textbf{Type} & \textbf{Key} & \textbf{Q} & \textbf{Type} & \textbf{Key} & \textbf{Q} & \textbf{Type} & \textbf{Key} & \textbf{Q} & \textbf{Type} & \textbf{Key} \\
\hline\hline
1 & MCQ & D & 15 & NAT & 2016.0 : 2016.0 & 29 & NAT & 0.33 : 0.34 & 43 & MCQ & B \\ \hline
2 & MCQ & A & 16 & MCQ & D & 30 & NAT & 28.0 : 28.0 & 44 & MCQ & C \\ \hline
3 & MCQ & A & 17 & MCQ & C & 31 & NAT & 456.0 : 456.0 & 45 & NAT & 9.0 : 9.0 \\ \hline
4 & MCQ & C & 18 & MCQ & B & 32 & NAT & 33.0 : 34.0 & 46 & MCQ & C \\ \hline
5 & MCQ & D & 19 & NAT & 10.0 : 10.0 & 33 & MCQ & B & 47 & NAT & 384.0 : 384.0 \\ \hline
6 & MCQ & B & 20 & MCQ & A & 34 & MCQ & D & 48 & NAT & 346.0 : 346.0 \\ \hline
7 & MCQ & D & 21 & MCQ & B & 35 & MCQ & A & 49 & NAT & 1.0 : 1.0 \\ \hline
8 & MCQ & D & 22 & MCQ & D & 36 & MCQ & D & 50 & MCQ & A \\ \hline
9 & MCQ & B & 23 & MCQ & B & 37 & MCQ & B & 51 & MCQ & A \\ \hline
10 & MCQ & B & 24 & MCQ & C & 38 & NAT & 12.0 : 12.0 & 52 & MCQ & B \\ \hline
11 & NAT & 6.0 : 6.0 & 25 & MCQ & C & 39 & NAT & 7.0 : 7.0 & 53 & NAT & 13.0 : 13.0 \\ \hline
12 & MCQ & D & 26 & NAT & 10.0 : 10.0 & 40 & MCQ & B & 54 & NAT & 1.1 : 1.1 \\ \hline
13 & MCQ & D & 27 & NAT & 197.9 : 198.1 & 41 & NAT & 256.0 : 256.0 & 55 & NAT & 2500.0 : 2500.0 \\ \hline
14 & MCQ & A & 28 & NAT & 2.0 : 2.0 & 42 & MCQ & D &  & &  \\ \hline
\end{tabular}%
}
\end{center}
\end{document}